% Annexe C -- Lean 4 pour debutants. A inclure apres \appendix.
% Les six exemples sont dans code/lean/appC_*.lean ; ils ont ete compiles le 2026-10-26 avec Lean 4.34.1 + Mathlib
% d13f23b7 (pile epinglee), et les sorties citees sont celles de ces compilations, copiees sans retouche.
% Definitions communes aux annexes C, D et E (chacune les porte, pour pouvoir etre incluse seule).
\ifdefined\chemin\else
\lstdefinestyle{appC}{language=Lean,
  literate={ℝ}{{$\mathbb{R}$}}1 {ℕ}{{$\mathbb{N}$}}1 {≤}{{$\le$}}1 {→}{{$\to$}}1 {∀}{{$\forall$}}1 {∃}{{$\exists$}}1
  {∧}{{$\wedge$}}1 {∈}{{$\in$}}1 {∑}{{$\sum$}}1 {·}{{$\cdot$}}1 {⟨}{{$\langle$}}1 {⟩}{{$\rangle$}}1 {α}{{$\alpha$}}1
  {é}{{\'e}}1 {è}{{\`e}}1 {ê}{{\^e}}1 {à}{{\`a}}1 {ô}{{\^o}}1 {ç}{{\c{c}}}1 {«}{{\guillemotleft}}1 {»}{{\guillemotright}}1
  {’}{{'}}1}
\DeclareUrlCommand\chemin{\urlstyle{tt}\def\UrlBreaks{\do\/\do\_\do\-\do\.}}
\DeclareUrlCommand\sha{\urlstyle{tt}\def\UrlBreaks{\do0\do1\do2\do3\do4\do5\do6\do7\do8\do9\do a\do b\do c\do d\do e\do f}}
\fi
\chapter{Lean 4 pour débutants}\label{appC:chap}

Le §\ref{ch1:sec:lean} a présenté Lean en lecteur : comment lire un énoncé et ce que garantit le noyau. Cette annexe
s'adresse à qui veut \emph{écrire} ses premières preuves. Elle suit six petits fichiers, rangés dans
\chemin{code/lean/appC_1_logique.lean} à \chemin{appC_6_contrefacon.lean}. Ils ont tous été compilés pour de vrai,
le 2026-10-26, avec la chaîne épinglée du livre (Lean~4.34.1, Mathlib au commit
\sha{d13f23b723b8a846827a245b89c10fc7d3f11612}). Chaque sortie reproduite ici est copiée de la compilation, sans
retouche, à ceci près que les chemins absolus des fichiers ont été raccourcis.

%=====================================================================================================================
\section{Installer Lean, Lake et Mathlib}\label{appC:sec:install}

\paragraph{Les trois outils.}
\begin{itemize}
\item \textbf{elan} est le gestionnaire de versions : il installe les chaînes Lean (\emph{toolchains}) et choisit la
bonne pour chaque projet, d'après le fichier \texttt{lean-toolchain} du projet.
\item \textbf{lake} est l'outil de construction, livré avec Lean. Il lit \texttt{lakefile.lean} (ou
\texttt{lakefile.toml}), télécharge les dépendances et compile les modules dans l'ordre de leurs \texttt{import}.
\item \textbf{Mathlib} est la bibliothèque mathématique communautaire~\cite{mathlib}. On en dépend comme d'un paquet
Lake, épinglé à un commit précis.
\end{itemize}

\paragraph{Installation type (Linux ou macOS).} Les commandes suivantes sont la procédure documentée par les projets
elan et Mathlib~\cite{app:elan,app:mathlibinstall} ; nous ne les avons \emph{pas} rejouées sur une machine vierge pour
ce livre (la machine de l'auteur avait déjà ses chaînes), et une version plus récente de ces outils peut en changer
les détails.
\begin{lstlisting}[style=appC,language={}]
curl https://raw.githubusercontent.com/leanprover/elan/master/elan-init.sh -sSf | sh
elan toolchain install leanprover/lean4:v4.34.1
lake new mon_projet math          # gabarit « math » : projet qui dépend de Mathlib
cd mon_projet
lake exe cache get                # télécharge les .olean précompilés de Mathlib
lake build
\end{lstlisting}
Pour reproduire exactement la version du livre, épinglez Mathlib dans \texttt{lakefile.lean} :
\begin{lstlisting}[style=appC,language={}]
require mathlib from git "https://github.com/leanprover-community/mathlib4.git" @ "d13f23b723b8a846827a245b89c10fc7d3f11612"
\end{lstlisting}
et mettez \chemin{leanprover/lean4:v4.34.1} dans \texttt{lean-toolchain}. Le cache précompilé n'existe que pour les
commits de Mathlib publiés ; sinon, \texttt{lake build} recompile Mathlib, ce qui prend des heures.

\paragraph{La chaîne du livre.} Les résultats Lean des chapitres~1 à~5 ont été vérifiés dans la copie du répertoire
\texttt{lean/} d'openai/math (avec sa Mathlib et ses quelque trente dépendances), et non dans un projet neuf
(annexe~\ref{appE:chap}, §\ref{appE:sec:lean}). Les six exemples de cette annexe n'importent que des modules de Mathlib
et du cœur de Lean ; un projet \texttt{lake new ... math} à la même révision de Mathlib devrait donc les compiler aussi,
mais nous ne l'avons pas vérifié.

\paragraph{Compiler un fichier.} Dans un projet, \texttt{lake env lean Fichier.lean} compile un fichier avec les
bonnes bibliothèques ; dans un éditeur (VS Code et l'extension Lean~4), l'état de la preuve s'affiche en continu.
Pour ce livre, nous avons utilisé le script du dépôt :
\begin{lstlisting}[style=appC,language={}]
taskset -c 6,7 .venv/bin/python -I scripts/certified_numerics/lean_pinned.py compile FICHIER.lean --json-out SORTIE.json
\end{lstlisting}
Il appelle \texttt{lean} \chemin{-DautoImplicit=false} avec le \texttt{LEAN\_PATH} de la pile épinglée et écrit le code de
retour, la durée et toute la sortie dans \texttt{SORTIE.json}.

\begin{warnbox}[: pas d'\texttt{import Mathlib} global]
Importez des modules précis (\lean{Mathlib.Tactic.Ring}, \lean{Mathlib.Tactic.Linarith}, ...), comme le font nos
exemples. Un \texttt{import Mathlib} charge toute la bibliothèque : c'est lent, et, dans une construction partielle de
Mathlib (comme celle du répertoire \texttt{formal/} du dépôt AutoevolveAI), cela échoue sur les modules non construits.
\end{warnbox}

%=====================================================================================================================
\section{Premiers théorèmes : logique et quantificateurs}\label{appC:sec:logique}

Le premier fichier n'importe rien : il n'utilise que le cœur de Lean.
\lstinputlisting[style=appC,title={\texttt{\detokenize{code/lean/appC_1_logique.lean}}}]{code/lean/appC_1_logique.lean}

\noindent Sortie de la compilation (code de retour $0$, 2,0~s) :
\begin{lstlisting}[style=appC,language={}]
'et_comm_terme' does not depend on any axioms
'et_comm_tactique' does not depend on any axioms
'modus_ponens_chaine' does not depend on any axioms
'existe_successeur' does not depend on any axioms
\end{lstlisting}

\paragraph{Ce qu'il faut retenir.}
\begin{itemize}
\item Après \lean{:=}, on donne soit un \emph{terme} (\texttt{$\langle$h.2, h.1$\rangle$}), soit \lean{by} suivi de \emph{tactiques}.
\item \lean{intro h} transforme le but $P\to Q$ en le but $Q$, avec l'hypothèse \lean{h : P} dans le contexte.
\item \lean{exact e} ferme le but si \lean{e} a exactement le type du but.
\item \lean{apply f} raisonne à rebours : si \texttt{f : A $\to$ B} et que le but est \lean{B}, le nouveau but est \lean{A}.
\item \lean{constructor} coupe un but $P\wedge Q$ en deux buts ; les points \texttt{$\cdot$} les traitent un par un.
\item \texttt{obtain $\langle$hp, hq$\rangle$ := h} découpe une hypothèse \texttt{h : P $\wedge$ Q} en ses deux morceaux.
\item Un « il existe » se prouve par une paire $\langle$témoin, preuve$\rangle$ ; ici, le témoin est $n+1$ et la preuve
\lean{rfl} (« égal par définition »).
\end{itemize}
Aucun de ces quatre théorèmes n'utilise d'axiome : la logique intuitionniste suffit.

%=====================================================================================================================
\section{Les tactiques usuelles}\label{appC:sec:tactiques}

\begin{table}[ht]\centering\small
\caption{Les tactiques employées dans les exemples, et ce qu'elles font.}\label{appC:tab:tactiques}
\begin{tabular}{p{2.4cm}p{11.2cm}}\toprule
tactique & effet\\\midrule
\lean{intro} & introduit l'hypothèse d'une implication ou la variable d'un $\forall$\\
\lean{exact} & ferme le but avec un terme du bon type\\
\lean{apply} & ramène le but à l'hypothèse d'une implication connue\\
\lean{constructor} & sépare un $\wedge$ (ou un $\leftrightarrow$) en deux buts\\
\lean{obtain}, \lean{rcases} & déstructure une hypothèse ($\wedge$, $\exists$, $\vee$)\\
\lean{rfl} & ferme un but $a=a$ « par définition »\\
\lean{simp} & réécrit avec les lemmes marqués \lean{@[simp]} (et ceux qu'on lui donne) jusqu'à une forme normale\\
\lean{ring} & prouve une égalité dans un (semi-)anneau commutatif (identité de polynômes)\\
\lean{linarith} & prouve une inégalité linéaire combinaison des hypothèses (ordre total, corps ordonné)\\
\lean{nlinarith} & comme \lean{linarith}, en multipliant des hypothèses ; on l'aide en lui fournissant des carrés\\
\lean{norm_num} & calcule exactement sur des littéraux numériques (entiers, rationnels)\\
\lean{positivity} & prouve $0<e$ ou $0\le e$ d'après la forme syntaxique de $e$\\
\lean{decide} & évalue une proposition décidable close (le noyau fait le calcul)\\
\lean{omega} & arithmétique linéaire sur $\N$ et $\Z$\\
\lean{induction} & récurrence sur un entier naturel (ou toute structure inductive)\\
\bottomrule\end{tabular}
\end{table}

\subsection{Calcul dans les réels}
\lstinputlisting[style=appC,title={\texttt{\detokenize{code/lean/appC_2_calcul.lean}}}]{code/lean/appC_2_calcul.lean}
\noindent Sortie (code de retour $0$, 6,3~s) :
\begin{lstlisting}[style=appC,language={}]
'carre_somme' depends on axioms: [propext, Classical.choice, Quot.sound]
'moyenne_entre' depends on axioms: [propext, Classical.choice, Quot.sound]
'am_gm_deux' depends on axioms: [propext, Classical.choice, Quot.sound]
'quatre_septiemes' depends on axioms: [propext, Classical.choice, Quot.sound]
'positif' depends on axioms: [propext, Classical.choice, Quot.sound]
\end{lstlisting}
Les nombres réels de Mathlib sont construits à partir des suites de Cauchy, avec l'axiome du choix : en pratique, tout théorème sur
$\R$ dépend des trois axiomes de la liste blanche. C'est normal.

\begin{remark}[Une bibliothèque qui bouge]
La première version de ce fichier importait \lean{Mathlib.Data.Real.Basic}. Elle compilait, avec cet avertissement :
\begin{lstlisting}[style=appC,language={}]
warning: 'Mathlib.Data.Real.Basic' has been deprecated: please replace this import by
import Mathlib.Basic.Real.Basic
\end{lstlisting}
Nous avons suivi le conseil. Les noms de modules et de lemmes changent d'une version de Mathlib à l'autre : c'est
pourquoi on épingle un commit, et pourquoi un fichier écrit pour une autre version peut ne plus compiler.
\end{remark}

\subsection{Récurrence, \texttt{simp} et \texttt{decide}}
\lstinputlisting[style=appC,title={\texttt{\detokenize{code/lean/appC_3_recurrence.lean}}}]{code/lean/appC_3_recurrence.lean}
\noindent Sortie (code de retour $0$, 4,3~s) :
\begin{lstlisting}[style=appC,language={}]
'double_somme' depends on axioms: [propext, Classical.choice, Quot.sound]
'petit_calcul' depends on axioms: [propext, Classical.choice, Quot.sound]
'simp_exemple' depends on axioms: [propext]
\end{lstlisting}
Dans \lean{double_somme}, le cas \lean{succ k} commence par \lean{rw [sum_range_succ]}, qui sépare le dernier terme :
$\sum_{i<k+1}i=\sum_{i<k}i+k$. Après normalisation par \lean{ring_nf}, l'hypothèse de récurrence \lean{ih} et le but sont
deux égalités d'entiers que \lean{omega} relie. Il y a souvent plusieurs chemins ; celui-ci n'est pas le plus court.

%=====================================================================================================================
\section{Lire un énoncé de Mathlib}\label{appC:sec:lire}

Avant d'utiliser un lemme, on lit son énoncé exact avec \texttt{\#check}. Le préfixe \lean{@} affiche aussi les arguments
implicites.
\lstinputlisting[style=appC,title={\texttt{\detokenize{code/lean/appC_4_lire_mathlib.lean}}}]{code/lean/appC_4_lire_mathlib.lean}
\noindent Sortie (code de retour $0$, 5,3~s) :
\begin{lstlisting}[style=appC,language={}]
Real.exp_pos : ∀ (x : ℝ), 0 < Real.exp x
Real.add_one_le_exp : ∀ (x : ℝ), x + 1 ≤ Real.exp x
@abs_sub_comm : ∀ {α : Type u_1} [inst : Lattice α] [inst_1 : AddGroup α] (a b : α), |a - b| = |b - a|
'exp_moins_positif' depends on axioms: [propext, Classical.choice, Quot.sound]
'exp_minore' depends on axioms: [propext, Classical.choice, Quot.sound]
\end{lstlisting}

\paragraph{Comment lire la troisième ligne.}
\begin{itemize}
\item \texttt{\{$\alpha$ : Type u\_1\}} entre accolades : un argument \emph{implicite}, que Lean devine (ici le type des éléments).
\item \texttt{[inst : Lattice $\alpha$]} entre crochets : une \emph{instance}, une structure que Lean cherche tout seul
(ici : $\alpha$ est un treillis, ce qui donne $\max$ et $\min$, donc la valeur absolue $|a|=\max(a,-a)$).
\item \texttt{(a b : $\alpha$)} entre parenthèses : les arguments explicites, qu'on donne à l'appel.
\item Après la virgule, la conclusion.
\end{itemize}
L'énoncé est donc : dans tout groupe additif réticulé, $|a-b|=|b-a|$. On l'applique comme une fonction : le terme
\lean{abs_sub_comm x y} a le type \lean{|x - y| = |y - x|}. C'est exactement ce que fait \lean{exp_moins_positif} avec
\lean{Real.exp_pos (-x)}. Les énoncés du livre se lisent de la même façon ; par exemple, dans
\texttt{triangular\_yukawa\_optimal \{$\kappa$ : $\mathbb R$\} (h$\kappa$ : 0 < $\kappa$) (C : Set Plane) ...} (chapitre~5), \texttt{$\kappa$} est implicite parce
qu'il se lit dans le type de \texttt{h$\kappa$}.

\paragraph{Chercher un lemme.} Les noms de Mathlib suivent une convention : \lean{add_one_le_exp} se lit
« $x+1\le\exp x$ ». La tactique \lean{exact?} cherche un lemme qui ferme le but, et le site de documentation de
Mathlib permet une recherche par nom. Le serveur LeanMaster utilisé pendant la rédaction offre aussi une recherche
(\texttt{search\_theorems}), mais sur sa propre bibliothèque.

%=====================================================================================================================
\section{Le contrôle des axiomes}\label{appC:sec:axiomes}

Un fichier qui compile n'est pas un fichier démontré. Le cinquième exemple contient deux faux théorèmes :
\lstinputlisting[style=appC,title={\texttt{\detokenize{code/lean/appC_5_temoins.lean}}}]{code/lean/appC_5_temoins.lean}
\noindent Sortie (code de retour $0$ !, 1,0~s) :
\begin{lstlisting}[style=appC,language={}]
appC_5_temoins.lean:8:8: warning: declaration uses `sorry`
'vrai_controle_positif' does not depend on any axioms
'faux_par_sorry' depends on axioms: [sorryAx]
'faux_par_axiome' depends on axioms: [triche]
\end{lstlisting}
Le compilateur accepte le fichier : \lean{sorry} ne produit qu'un avertissement, et un \lean{axiom} est une déclaration
légale. Ce sont les lignes \texttt{\#print axioms} qui révèlent la triche : \lean{sorryAx} dans la première, l'axiome
utilisateur \lean{triche} dans la seconde. La règle de ce livre est donc :
\begin{center}
\emph{un théorème n'est retenu que si \texttt{\#print axioms} n'affiche que des axiomes de
$\{$\lean{propext}, \lean{Classical.choice}, \lean{Quot.sound}$\}$.}
\end{center}
Le théorème \lean{vrai_controle_positif} est le témoin positif : la porte le laisse passer.

%=====================================================================================================================
\section{Pourquoi il faut aussi rejouer le noyau}\label{appC:sec:leanchecker}

\subsection{Une contrefaçon invisible à \texttt{\#print axioms}}
Le noyau de Lean est le seul juge, mais le compilateur peut lui demander de ne pas juger : l'option
\lean{debug.skipKernelTC} désactive la vérification de types du noyau pour les déclarations ajoutées. Le sixième
fichier s'en sert pour fabriquer une « preuve » de \lean{False} :
\lstinputlisting[style=appC,title={\texttt{\detokenize{code/lean/appC_6_contrefacon.lean}}}]{code/lean/appC_6_contrefacon.lean}
\noindent Sortie (code de retour $0$ ; 51,6~s à la première compilation, cache disque froid, 4,4~s à la seconde) :
\begin{lstlisting}[style=appC,language={}]
'contrefacon' does not depend on any axioms
'un_egale_deux' does not depend on any axioms
\end{lstlisting}
Le fichier compile, et \texttt{\#print axioms} déclare que $1=2$ ne dépend d'\emph{aucun} axiome. La porte de la
section précédente le laisserait passer.

\begin{remark}[Un premier essai, refusé]
Notre première version écrivait simplement \lean{set_option debug.skipKernelTC true in theorem contrefacon : False :=
True.intro}. Elle a été refusée (code de retour $1$) avec le message \texttt{Type mismatch: True.intro has type True
but is expected to have type False} : avant le noyau, l'\emph{élaborateur} vérifie lui aussi les types. L'option ne
débranche que le noyau ; pour la contrefaçon, il faut donc ajouter la déclaration directement au noyau
(\lean{addDecl}), comme l'a fait l'auditeur de N4. Cette précision ne rassure pas : un fichier hostile peut faire
exactement cela.
\end{remark}

\subsection{Le rejeu par \texttt{leanchecker}}
L'outil \lean{leanchecker}, livré avec Lean, relit le fichier \texttt{.olean} compilé et \emph{rejoue} chaque
déclaration du module dans un noyau neuf, sans faire confiance à ce que le compilateur a enregistré. Il faut donc
d'abord produire un \texttt{.olean}. Le script \chemin{code/appC_leanchecker.py} compile chaque exemple en module
\lean{AppC.E<n>_<nom>} dans un dossier de travail, puis lance \lean{leanchecker} avec le même \texttt{LEAN\_PATH} :
\begin{lstlisting}[style=appC,language={}]
taskset -c 6,7 .venv/bin/python -I papers/textbook_formalisation_2026/code/appC_leanchecker.py \
    DOSSIER_DE_TRAVAIL SORTIE.json appC_1_logique appC_2_calcul appC_3_recurrence \
    appC_4_lire_mathlib appC_5_temoins appC_6_contrefacon
\end{lstlisting}
Résultat du 2026-10-26, pour les six exemples puis un module inexistant (compilation et rejeu, deux c\oe urs) :
\begin{center}\small
\begin{tabular}{lrrrl}\toprule
module & compil. & rejeu & durée & sortie du rejeu\\\midrule
\texttt{AppC.E1\_logique} & 0 & 0 & 48,4~s & (vide)\\
\texttt{AppC.E2\_calcul} & 0 & 0 & 54,3~s & (vide)\\
\chemin{AppC.E3_recurrence} & 0 & 0 & 56,7~s & (vide)\\
\chemin{AppC.E4_lire_mathlib} & 0 & 0 & 58,4~s & (vide)\\
\texttt{AppC.E5\_temoins} & 0 & 0 & 53,6~s & (vide)\\
\chemin{AppC.E6_contrefacon} & 0 & \textbf{1} & 58,6~s & voir ci-dessous\\
\chemin{AppC.ModuleInexistant} & -- & \textbf{1} & -- & \texttt{Could not find any oleans}\\
\bottomrule\end{tabular}
\end{center}
La sortie du rejeu de l'exemple~6 est :
\begin{lstlisting}[style=appC,language={}]
leanchecker found a problem in AppC.E6_contrefacon
uncaught exception: while replaying declaration 'contrefacon':
(kernel) declaration type mismatch, 'contrefacon' has type
  True
but it is expected to have type
  False
\end{lstlisting}
Le script se termine par la ligne \texttt{\{"all\_controls\_pass": true\}} : chaque module a rendu le code attendu.

\paragraph{Lecture.}
\begin{itemize}
\item \emph{Témoins positifs.} Les modules honnêtes sont rejoués sans erreur (code~$0$, aucune sortie).
\item \emph{Témoin négatif principal.} La contrefaçon est refusée (code~$1$) : en rejouant \lean{contrefacon}, le noyau
constate que le terme \lean{True.intro} a le type \lean{True} et non \lean{False}.
\item \emph{Témoin négatif trivial.} Un module inexistant est refusé (code~$1$) : le code~$0$ des autres n'est pas
celui d'un outil qui ne fait rien.
\item Une minute environ par module, même pour les exemples~1 et~5 qui n'importent rien : le temps est celui du démarrage de
\lean{leanchecker}, pas celui du module (nous n'en avons pas mesuré la cause). Le rejeu ne porte que sur les déclarations du module : Mathlib et le cœur de Lean sont chargés,
pas rejoués (il faudrait l'option \texttt{--fresh}, beaucoup plus coûteuse).
\end{itemize}

\begin{warnbox}[: chaque contrôle attrape ce que l'autre laisse passer]
\lean{leanchecker} ne remplace pas \texttt{\#print axioms} : un \lean{sorry} ou un \lean{axiom} déclaré est
\emph{légal} pour le noyau, et le rejeu du module de l'exemple~5 est accepté (tableau ci-dessus : code~$0$ ; c'est le comportement attendu). Inversement, \texttt{\#print axioms} ne voit pas la contrefaçon de l'exemple~6.
D'où la porte de niveau~A du livre, en quatre contrôles (§\ref{ch1:sec:lean}) :
\begin{enumerate}
\item la compilation sans erreur ;
\item la liste des axiomes incluse dans la liste blanche ;
\item le rejeu \lean{leanchecker} du module, avec le code~$0$ ;
\item une recherche textuelle de \lean{skipKernelTC}, \lean{addDecl}, \lean{native_decide} et \lean{implemented_by}
dans les sources (elle trouverait l'exemple~6 immédiatement).
\end{enumerate}
Quand le troisième contrôle manque, le livre écrit \tierAs{} au lieu de \tierA.
\end{warnbox}

\paragraph{Et \texttt{native\_decide} ?} La tactique \lean{native_decide} fait évaluer une proposition par du code
compilé, hors du noyau. Lean~4.34.1 l'enregistre comme un axiome propre à chaque usage (de la forme
\texttt{\_native.native\_decide.ax\_\dots}), que \texttt{\#print axioms} affiche et que la liste blanche refuse ; LeanFlow
le classe à part (annexe~\ref{appD:chap}). Pour faire calculer le noyau lui-même, on écrit \lean{decide +kernel}.

%=====================================================================================================================
\section{Exercices}\label{appC:sec:exos}

\begin{exercise}
Écrire et compiler la preuve de $\forall\,p\,q:\mathrm{Prop},\ (p\vee q)\to(q\vee p)$, d'abord par un terme, puis avec
les tactiques \lean{intro}, \lean{rcases} et \lean{exact} (indications : \lean{Or.inl}, \lean{Or.inr}).
\end{exercise}

\begin{exercise}
Prouver en Lean que $\sum_{i<n}(2i+1)=n^2$ sur le modèle de \lean{double_somme}. Vérifier avec \texttt{\#print axioms}
que votre preuve n'utilise que la liste blanche.
\end{exercise}

\begin{exercise}
Dans \chemin{appC_2_calcul.lean}, remplacer \lean{nlinarith [sq_nonneg (a - b)]} par \lean{linarith}. Que se
passe-t-il, et pourquoi ?
\end{exercise}

\begin{exercise}[Témoin négatif]
Modifier \lean{quatre_septiemes} en $4/7-1/3=1/4$ et compiler. Quel est le message, et quel est le code de retour ?
Comparer avec l'exemple~5 : pourquoi une preuve fausse fait-elle échouer la compilation, alors qu'un \lean{sorry} ne la
fait pas échouer ?
\end{exercise}
